\documentclass[12pt,a4paper]{article}
\usepackage[utf8]{inputenc}
\usepackage[T1]{fontenc}
\usepackage[brazilian]{babel}
\usepackage{amsmath,amssymb}
\usepackage[margin=2.2cm]{geometry}
\usepackage{tikz}
\usetikzlibrary{arrows.meta,decorations.pathmorphing,positioning}
\usepackage{tcolorbox}
\tcbuselibrary{skins,breakable}
\usepackage{enumitem}
\usepackage{xcolor}
\usepackage{hyperref}
\hypersetup{colorlinks=true,linkcolor=blue!50!black,urlcolor=blue!50!black}

\newtcolorbox{definicao}[1]{colback=blue!5,colframe=blue!55!black,fonttitle=\bfseries,coltitle=white,title={#1},breakable,boxrule=0.8pt,arc=2mm}
\newtcolorbox{exemplo}[1]{colback=green!5,colframe=green!35!black,fonttitle=\bfseries,coltitle=white,title={#1},breakable,boxrule=0.8pt,arc=2mm}
\newtcolorbox{atencao}[1]{colback=red!5,colframe=red!55!black,fonttitle=\bfseries,coltitle=white,title={#1},breakable,boxrule=0.8pt,arc=2mm}
\newtcolorbox{dica}[1]{colback=yellow!12,colframe=orange!70!black,fonttitle=\bfseries,coltitle=black,title={#1},breakable,boxrule=0.8pt,arc=2mm}

% --- Ferramenta para desenhar retas de intervalo ---
% \retaIntervalo{esquerda}{direita}{tipo esquerdo}{tipo direito}{extra}
% tipo: "fechado" (bola cheia) ou "aberto" (bola vazia)
\newcommand{\bolaFechada}[1]{\filldraw[blue] (#1,0) circle (2.2pt);}
\newcommand{\bolaAberta}[1]{\draw[blue,fill=white,thick] (#1,0) circle (2.2pt);}

\title{\vspace{-1.5cm}\textbf{Tema 2 — Intervalos Reais}\\[2pt]\large A "tradução" entre desigualdades, reta numérica e notação de colchetes}
\author{}
\date{}

\begin{document}
\maketitle
\thispagestyle{empty}

\section*{Por que isso importa?}
Toda vez que você resolve uma inequação, o resultado é um \textbf{intervalo}
— um pedaço da reta real. E todo domínio de função (Tema 6) é escrito como
intervalo. Dominar as 3 formas de representar a mesma informação
(desigualdade $\to$ reta $\to$ colchetes) é obrigatório.

\section{As 4 formas de um intervalo}

\begin{definicao}{Intervalo fechado $[a,b]$}
Inclui os extremos $a$ e $b$.
\[
[a,b] = \{x \in \mathbb{R} \mid a \le x \le b\}
\]
\begin{center}
\begin{tikzpicture}[scale=1]
  \draw[-{Stealth[length=2.5mm]}, thick] (-1,0) -- (5,0);
  \foreach \x in {0,1,2,3,4} \draw (\x,0.06)--(\x,-0.06);
  \node[below] at (0,-0.15) {\small $a$};
  \node[below] at (4,-0.15) {\small $b$};
  \draw[ultra thick, blue] (0,0) -- (4,0);
  \bolaFechada{0} \bolaFechada{4}
\end{tikzpicture}
\end{center}
\end{definicao}

\begin{definicao}{Intervalo aberto $(a,b)$ ou $\,]a,b[\,$}
\textbf{Exclui} os extremos.
\[
(a,b) = \{x \in \mathbb{R} \mid a < x < b\}
\]
\begin{center}
\begin{tikzpicture}[scale=1]
  \draw[-{Stealth[length=2.5mm]}, thick] (-1,0) -- (5,0);
  \foreach \x in {0,1,2,3,4} \draw (\x,0.06)--(\x,-0.06);
  \node[below] at (0,-0.15) {\small $a$};
  \node[below] at (4,-0.15) {\small $b$};
  \draw[ultra thick, blue] (0,0) -- (4,0);
  \bolaAberta{0} \bolaAberta{4}
\end{tikzpicture}
\end{center}
\end{definicao}

\begin{definicao}{Intervalos mistos (semiabertos)}
\[
[a,b) = \{x \mid a \le x < b\} \qquad\qquad (a,b] = \{x \mid a < x \le b\}
\]
\begin{center}
\begin{tikzpicture}[scale=1]
  \draw[-{Stealth[length=2.5mm]}, thick] (-1,0) -- (5,0);
  \foreach \x in {0,1,2,3,4} \draw (\x,0.06)--(\x,-0.06);
  \node[below] at (0,-0.15) {\small $a$};
  \node[below] at (4,-0.15) {\small $b$};
  \draw[ultra thick, blue] (0,0) -- (4,0);
  \bolaFechada{0} \bolaAberta{4}
\end{tikzpicture}
\end{center}
\end{definicao}

\begin{dica}{Como lembrar bola cheia $\times$ bola vazia}
\textbf{Bola cheia (fechada)} = ``o número entra'' = símbolo $\le$ ou $\ge$
= colchete \textbf{para dentro} \texttt{[ ]}. \textbf{Bola vazia (aberta)} =
``o número não entra'' = símbolo $<$ ou $>$ = colchete \textbf{para fora}
\texttt{] [} ou parêntese \texttt{( )}. As duas notações $[a,b[$ e $[a,b)$
significam exatamente a mesma coisa — a primeira é mais comum no Brasil, a
segunda é mais internacional.
\end{dica}

\section{Intervalos infinitos}

Quando o intervalo não tem um dos lados limitado, usamos $+\infty$ ou
$-\infty$ — e o lado do infinito \textbf{sempre} é aberto (infinito não é um
número, então nunca ``inclui'' o infinito).

\begin{exemplo}{$x \ge 3$}
\[
[3,+\infty[ \ = \{x\in\mathbb{R}\mid x\ge 3\}
\]
\begin{center}
\begin{tikzpicture}
  \draw[-{Stealth[length=2.5mm]}, thick] (-1,0) -- (5,0);
  \draw (3,0.06)--(3,-0.06) node[below]{\small $3$};
  \draw[ultra thick, blue, -{Stealth[length=3mm]}] (3,0) -- (4.8,0);
  \bolaFechada{3}
\end{tikzpicture}
\end{center}
\end{exemplo}

\begin{exemplo}{$x \le 4$}
\[
\,]-\infty,4]\ = \{x\in\mathbb{R}\mid x\le 4\}
\]
\begin{center}
\begin{tikzpicture}
  \draw[-{Stealth[length=2.5mm]}, thick] (-1,0) -- (5,0);
  \draw (4,0.06)--(4,-0.06) node[below]{\small $4$};
  \draw[ultra thick, blue, -{Stealth[length=3mm]}] (4,0) -- (-0.8,0);
  \bolaFechada{4}
\end{tikzpicture}
\end{center}
\end{exemplo}

\section{Do desenho para a notação (e vice-versa)}

\begin{exemplo}{Lendo um desenho com dois pontos, ambos abertos}
\begin{center}
\begin{tikzpicture}
  \draw[-{Stealth[length=2.5mm]}, thick] (-1,0) -- (5,0);
  \draw (1,0.06)--(1,-0.06) node[below]{\small $1$};
  \draw (4,0.06)--(4,-0.06) node[below]{\small $4$};
  \draw[ultra thick,blue,decorate,decoration={snake,amplitude=1pt,segment length=4pt}] (1,0) -- (4,0);
  \bolaAberta{1} \bolaFechada{4}
\end{tikzpicture}
\end{center}
A bolinha em $1$ é vazia (aberto) e em $4$ é cheia (fechado). Traduzindo:
\[
1 < x \le 4 \qquad\Longleftrightarrow\qquad \,]1,4]
\]
\end{exemplo}

\begin{atencao}{Rabiscos ``rasurados'' na reta}
Muitas listas (inclusive a sua) desenham um rabisco tipo mola
(\emph{ziguezague}) na parte da reta que \textbf{não} pertence ao
intervalo, e deixam \textbf{limpo} (ou tracejado) o trecho que pertence.
Preste atenção: o rabisco indica ``isso aqui está fora da solução'', não o
contrário!
\end{atencao}

\section{Operações entre intervalos: união e interseção}

Isso é essencial no Tema 6 (domínio de função), quando você precisa
combinar duas condições ao mesmo tempo (ex.\ denominador $\neq 0$
\textbf{e} raiz $\ge 0$).

\begin{definicao}{Interseção $A \cap B$ — ``e'' (ao mesmo tempo)}
É o que está em \textbf{ambos} os conjuntos simultaneamente.
\end{definicao}

\begin{definicao}{União $A \cup B$ — ``ou''}
É tudo que está em \textbf{pelo menos um} dos dois conjuntos.
\end{definicao}

\begin{exemplo}{$[1,7] \cap\, ]3,+\infty[$}
\begin{center}
\begin{tikzpicture}[yscale=0.8]
  \draw[-{Stealth[length=2.5mm]}, thick] (-1,0) -- (8,0);
  \foreach \x in {1,3,7} \draw (\x,0.06)--(\x,-0.06) node[below]{\small \x};

  \draw[ultra thick, blue] (1,0.35) -- (7,0.35);
  \bolaFechada{1} \node[blue] at (1,0.35+0.001){}; \filldraw[blue] (1,0.35) circle(2.2pt); \filldraw[blue] (7,0.35) circle(2.2pt);
  \node[left, blue] at (-1.1,0.35) {\small $A$};

  \draw[ultra thick, orange] (3,-0.35) -- (7.7,-0.35);
  \draw[orange,fill=white,thick] (3,-0.35) circle (2.2pt);
  \node[left, orange] at (-1.1,-0.35) {\small $B$};

  \draw[ultra thick, red] (3,0) -- (7,0);
  \draw[red,fill=white,thick] (3,0) circle(2.2pt); \filldraw[red] (7,0) circle(2.2pt);
  \node[left, red] at (-1.1,0) {\small $A\cap B$};
\end{tikzpicture}
\end{center}
Resultado: $\,]3,7]$ (aberto onde qualquer um dos dois for aberto naquele
ponto; sempre pegue o trecho mais ``restrito'').
\end{exemplo}

\section*{Resumo relâmpago}
\begin{center}
\renewcommand{\arraystretch}{1.4}
\begin{tabular}{lll}
\hline
\textbf{Desigualdade} & \textbf{Notação} & \textbf{Bolinha} \\
\hline
$a \le x \le b$ & $[a,b]$ & cheia -- cheia \\
$a < x < b$ & $\,]a,b[\,$ ou $(a,b)$ & vazia -- vazia \\
$a \le x < b$ & $[a,b[\,$ ou $[a,b)$ & cheia -- vazia \\
$x \ge a$ & $[a,+\infty[$ & cheia -- (seta) \\
$x < b$ & $\,]-\infty,b[$ & (seta) -- vazia \\
\hline
\end{tabular}
\end{center}

\end{document}
